% 本文件由 LaTeX 排版 Agent 根据有效 translation.json 自动生成。
% author=Tertullian; work_id=0405; title=论劝勉贞洁

\chapter{论劝勉贞洁}

% structure: p0001 source=h1 target=omitted reason=page-title-already-rendered-by-parent

% structure: p0002 source=h3 target=section reason=relative-heading-rank; base=h3

\section{第一章 引言。童贞分为三种}

\Emph{我不怀疑}，兄弟，在你的妻子平安离世后，你全心致力于调整心灵（到正确状态），正在认真思考独身生活的结局，当然需要指导。虽然在此类情形中，每个人都应当与自己的信德交谈，并咨询其力量；然而，由于在这（特定的）试炼类型中，\Emph{肉身的需要}（它通常是信德在同一内在良知的法庭上的对手，如我所述）激起了思虑，信德需要外来的指导，如同辩护者，来对抗\Emph{肉身的需要}：这种需要其实很容易被限制，如果考虑的是天主的\Emph{旨意}而非\Emph{宽仁}的话。没有人因利用宽仁而配得（恩宠），而是因对（主人的）旨意的迅速服从。天主的旨意是我们的圣化，因为祂愿意祂的\Quote{肖像}——我们——也成为祂的\Quote{模样}；使我们\Quote{圣洁}，正如祂自己\Quote{圣洁}。那善——即圣化——我将其分为几种类型，以便我们在某一种类型中被找到。第一种类型是生来童贞；第二种是从\Emph{第二次}出生而来的童贞，即从圣洗池而来；这种（第二次童贞）在婚姻状态中借着共同协议保持纯洁，或者自愿坚持守寡；还有第三个等级，即一夫一妻制，即在第一次婚姻中断后，从此放弃性关系。第一种童贞是幸福的（童贞），完全不知道你将后来希望从中解脱的事物；第二种是德性的（童贞），轻视你完全知晓其力量的事物；最后一种，即在婚姻因死亡而解除后不再婚，除了是德性的光荣外，也是节制的光荣；因为节制是不惋惜已被夺去的事物，且是被主天主夺去，若没有祂的旨意，就没有一片叶子从树上飘落，也没有一只价值一分钱的麻雀掉到地上。\par

% structure: p0004 source=h3 target=section reason=relative-heading-rank; base=h3

\section{第二章 我们的过犯不应归咎于天主。人唯一拥有的能力是意志的能力}

简而言之，在那句话中有什么节制呢：\Quote{主赐予，主又收回；按主的意愿成就了！}因此，如果我们重新缔结已被收回的婚姻，无疑是在对抗天主的旨意，愿意再次拥有祂不愿我们拥有的事物。因为如果祂愿意（我们拥有），就不会收回它；除非我们也将此解释为天主的旨意，仿佛祂再次愿意我们拥有祂刚才所不愿的。忠诚而坚定的信德不应以这种方式将一切归于天主的旨意，也不应让每个人如此自欺，说\Quote{没有祂的许可，什么事也不会发生}，以至于让我们无法理解有某事属于我们自己的能力。否则，如果我们坚持认为没有天主的旨意我们什么也做不了，那么每一种罪都将被原谅；这个定义将摧毁整个纪律，甚至毁灭天主本身；如果祂按其旨意产生祂所不愿的事物，或者根本没有天主所不愿的事物。但是，正如有些事物是祂所禁止的，祂甚至为此宣告永罚——当然，祂所\Emph{禁止}并因此\Emph{触怒}祂的事物，祂并不\Emph{愿意}——同样，相反地，祂\Emph{确实}愿意的，祂命令并视为可接纳的，并以永恒的赏报来回报。因此，当我们从祂的诫命中得知每一类行为，祂所不愿的和祂所愿意的，我们仍然拥有意愿和选择其一的能力；正如经上所记载的：\Quote{看，我已将善与恶摆在你面前，因为你尝了知善恶树的果子。}因此，我们不应将属于我们自己选择的事归于主的旨意；（假设）祂不愿意，或者（积极）不愿善，那不愿恶的（主）也是如此。所以，当我们愿意恶事，与愿意善的天主旨意相背时，这是出于我们自己的意愿。此外，若你问那使我们意愿任何与天主旨意相背之事的意愿从何而来，我要说，它源于我们自己。我说这话并非鲁莽——因为你必须与你所从出的种子相符——倘若确实如此，即我们族类与原罪的开创者亚当，他自己意愿了他所犯的罪。因为魔鬼并没有将犯罪的意愿强加于他，只是为那意愿提供了素材。另一方面，天主的旨意则成了服从的问题。同样，你，若你不服从那以自由行动的诫命训练你的天主，将通过你意志的自由，心甘情愿地转向做天主所不愿的事，并因此认为自己被魔鬼颠覆了；然而他虽然\Emph{愿意}你意愿天主所不愿的事物，却并未\Emph{使}你意愿，正如他并未\Emph{使}我们那些原祖沦入犯罪的意愿；不，也并未使他们违背自己的意愿，或在不知天主所不愿的情况下（迫使）。因为，当然，天主不愿某事被行时，就使死亡成为其必然后果。因此，魔鬼的作为只有一件：试探你是否愿意那取决于你意愿之事。但当你已经\Emph{意愿}时，他就因此把你置于自己的权下；并非由于他在你内\Emph{造成}了意愿，而是由于在你的意愿中找到了可乘之机。所以，既然唯一在我们能力之内的是意愿——而正是在此我们的心被考验，看我们是否愿意符合祂旨意的事物——我们必须深刻而焦虑地反复思考天主的旨意，我说，（要看）甚至祂可能愿意什么隐秘之事。\par

% structure: p0006 source=h3 target=section reason=relative-heading-rank; base=h3

\section{第三章 论宽容与纯粹意愿，并以此问题为例证。}

因为什么是显而易见的，我们都知道；而这些事物究竟在何种意义上是显而易见的，则必须彻底考查。因为，尽管有些事似乎带有\Quote{天主的旨意}的意味，既然它们被祂所允许，但并不能直接推论出每一件许可之事都出自允许者纯粹而绝对的旨意。宽容是所有许可的根源。尽管宽容并非独立于意愿，但既然它有其原因在于被宽容者，它（可以说）出自不情愿的意愿，经历了一个对其自身产生迫使意愿的成因。请看，一个以他人为原因的意愿其性质如何。此外，还有第二种纯粹意愿需要加以考虑。天主愿意我们做一些取悦于祂的行为，在这些行为中，并非宽容在支持，而是纪律在主导。然而，如果祂将某些别的行为置于这些之上——当然，这些是祂更愿意的行为——那么毫无疑问，我们应当追求的是祂更愿意的行为；因为那些祂较不愿意的（由于祂更愿意别的）应被视为如同祂不愿意一样。因为，通过显示祂更愿意什么，祂已以较大的意愿消除了较小的意愿。而祂向你的知识提出每一种（意愿）的程度，也正是祂定义你的责任去追求祂所宣告祂更愿意之事。那么，如果祂宣告的目的在于使你追求祂更愿意之事，无疑，除非你这样做，你就带有与祂意愿相背的意味，因为带有与祂最高意愿相背的意味；并且你与其说是赚取赏报，不如说是冒犯，因为你固然做了祂所愿意的事，却拒绝了祂更愿意的事。部分来说，你犯罪了；部分来说，即使你没有犯罪，你仍然不配得赏报。此外，连不愿配得赏报本身难道不就是一种罪吗？\par

因此，如果再婚在其许可的来源上，在于那被称为宽免的\Quote{天主的旨意}，我们将否认那以宽免为原因的是纯粹的意愿；如果其来源在于那种被另一事物——即那视节欲更为可取的事物——所偏好而成为更好的，我们将从上述论证中学到，不更优的被更优的废止。请允许我谈及这些考虑，以便我现在可以遵循宗徒言论的脉络。但首先，我若评述他自己所声称的，我不会被认为不虔敬；即他引入了所有关于婚姻的\Emph{宽免}，是凭他自己的判断——也就是出于人的感觉，而非出于天主的训令。因为，此外，当他针对\BibleQuote{寡妇和未婚者}明确规则时，即他们\BibleQuote{若不能自制，就应当结婚}，因为\BibleQuote{结婚比欲火攻心更好}，他转向另一类人说：\BibleQuote{至于已婚者，我正式宣告——不是我，而是主。}这样他显示，通过将他自身位格转给主，他上面所说的话并非以主的位格宣告，而是以他自己的位格：\BibleQuote{结婚比欲火攻心更好}。现在，虽然那句话涉及那些在未婚或寡居状态下被信德所\Emph{\Quote{抓住}}的人，然而，既然所有人都以此为借口，倾向于婚姻中的放纵，我愿意深入探讨他所指出的那种善，即\Quote{比惩罚更好}的善；这种善若非与极坏之物比较，便不能显得是善；因此\Quote{结婚}之所以是善，是因为\Quote{欲火攻心}更坏。\Quote{善}之所以配称为善，是在于它在不作比较时仍能保持此名；我所说的比较，不是与恶比，甚至不是与另一种次等善相比；因此，即使它\Emph{是}与某种其他善相比较，并被另一种善所遮蔽，它仍然保持着\Quote{善}的名号。然而，如果是一种\Emph{恶}的本性迫使我们将它称为\Quote{善}，那么它与其说是\Quote{善}，不如说是一种较低级的恶，被一种高级的恶所遮蔽，从而被迫冠以善之名。简言之，除去比较的条件，不要说\BibleQuote{结婚比欲火攻心更好}；我怀疑你是否有胆量只说\Quote{结婚更好}，而不补充说明什么更好。因此，不\Emph{更优}的当然也不\Emph{善}；因为你已除去比较的条件，而比较的条件在使事物成为\Quote{更好}的同时，也迫使它被视为\Quote{善}。\BibleQuote{结婚比欲火攻心更好}的理解与\Quote{少一只眼睛比两只眼睛都要好}相同：然而，如果你撤去比较，那么有一只眼睛就不再是\Quote{更好}，因为它也不是\Quote{善}。因此，不要有人从这段经文中抓取婚姻的辩护，这段经文原本是指\Quote{未婚者和寡妇}，对他们来说，还没有婚姻的结合；虽然我希望我已经表明，连这样的人也必须理解\Emph{许可}的本质。\par

% structure: p0009 source=h3 target=section reason=relative-heading-rank; base=h3

\section{第四章 进一步评述宗徒言论}

然而，关于第二次婚姻，我们清楚地知道宗徒已宣告：\BibleQuote{你已从妻子身上被解脱；不要寻求妻子。但如果你结婚，你并不会犯罪。} 然而，正如前一种情况，他也从个人的劝告、而非从天主的诫命引入了这一论述的顺序。但天主的诫命与人的劝告之间相差甚远。他说：\BibleQuote{主的诫命，我没有；但我给出忠告，如同从主那里获得怜悯而成为忠信的人。} 事实上，无论在福音还是在保禄自己的书信中，你都不会找到天主的诫命作为允许重复婚姻的来源。由此，婚姻的合一必须遵守的教义得到了确认；因为凡未发现为主所\Emph{允许}的，便应承认为\Emph{禁止}的。再加上这个事实：就连这种人的劝告的引入，仿佛已经开始反省自身的过分，立即约束并收回了自己，同时附加上：\BibleQuote{然而，这样的人将在肉体上遭遇压力}；当他说他\BibleQuote{宽免他们}时；当他补充说\BibleQuote{时期缩短了}，以致\BibleQuote{就连有妻子的，也当像没有一样}；当他比较已婚者和未婚者的挂虑时：因为他通过这些考虑教导结婚为何不适宜的理由，从而劝阻了他在前面曾予以宽容的事。对于第一次婚姻尚且如此，何况对于第二次呢！然而，当他劝勉我们效法他自己的榜样时，当然是在表明他\Emph{确实}希望我们成为什么，即是节制的；他同样宣告了他\Emph{不}希望我们成为什么，即是\Emph{不}节制的。同样，当他\Emph{意愿}一件事时，他并不自发地或真正地允许他所不愿的事。因为他如果愿意，就不会\Emph{允许}；反而会\Emph{命令}。\BibleQuote{但是再看：女人若丈夫死了，他说，便可结婚，只要她愿意嫁给任何人，只要\InnerQuote{在主内}。} 啊！但他说\BibleQuote{她将更幸福}，\BibleQuote{如果她照我的意见，永久保持现状。而且，我认为，我也拥有圣神。} 我们看见两种劝告：一种是他上文给予结婚的宽容；另一种是他随后教导关于结婚的节制。你说：\Quote{那么，我们该同意哪一个呢？} 仔细看看，然后选择。在给予宽容时，他称引一个明智\Emph{人}的劝告；在命令节制时，他肯定\Emph{圣神}的劝告。要遵循那以神性为保护者的训诫。的确，信徒们同样\BibleQuote{拥有圣神}；但并非所有信徒都是宗徒。因此，当那个曾自称\Quote{信徒}的人，随后又说他\BibleQuote{拥有圣神}——即使在普通信徒身上，也无人会怀疑这一点——他这样说，是为了重申自己的宗徒尊严。因为宗徒们以本有的方式拥有圣神，他们完满地拥有祂，在先知预言、异能效力、以及舌音证据的运作中；而不像其他人那样只局部拥有。因此，他将圣神的权威附加于他宁愿我们留意的那一种劝告形式上；于是它立即不再是圣神的\Emph{劝告}，而考虑到祂的尊威，成了\Emph{诫命}。\par

% structure: p0011 source=h3 target=section reason=relative-heading-rank; base=h3

\section{第五章 婚姻的合一由婚姻的最初建立以及宗徒将那原初模型应用于基督与教会所教导}

因为确立一次结婚的法律，人类本身的起源就是我们的权威；它有力地见证天主在起初制定了什么作为后世要仔细考察的模型。因为当祂塑造了人，并预见到了需要一个与他相配的伴侣时，祂从他的肋骨中取出一根，为他造了一个女人；然而，当然，创造主或材料都不会不足以（创造更多）。亚当有更多的肋骨，天主有不知疲倦的手；但在天主的眼中并没有更多的妻子。因此，天主的人亚当和天主的女人厄娃，彼此履行一次婚姻的责任，藉着他们起源的权威先例和天主的原始旨意，为人类确立了一个模型。最后，祂说：\BibleQuote{将来，两个人成为一体，不是三个也不是四个}。在任何其他假设下，就不再是\Quote{一体}，也不是\Quote{两个人（结合）成为一体}。如果结合和合而为一是一次性的，就会如此。但如果第二次或更频繁地发生，肉体立刻就不再是\Quote{一体}，也不再是\Quote{两个人（结合）成为一体}，而显然是一根肋骨（分开）成更多的。但是当宗徒将\BibleQuote{两个人将成为一体}解释为指向教会与基督，按照教会与基督的属灵婚配（因为基督是唯一的，祂的教会也是唯一的），我们必须认识到婚姻合一的律法对我们有双重的和额外的强调，不仅按照我们族类的根基，也按照基督的圣事。在这两种情况下，我们都从一次婚姻中获得起源；在亚当内属血肉，在基督内属神灵。两次诞生共同确立了一夫一妻制的规范。就这两者而言，凡越过一夫一妻界限的，都是堕落的。多妻制始于一个被诅咒的人。拉默客是第一个娶两个女人为妻的人，使\Emph{三个人}（结合）成为\Quote{一体}。\par

% structure: p0013 source=h3 target=section reason=relative-heading-rank; base=h3

\section{第六章 对族长多妻制的异议之回应}

\Quote{然而那些蒙福的圣祖们}，你说，\Quote{不仅与多位妻子，也与妾室通婚。} 那么，这会使我们也可以无限制地结婚吗？我承认它会，如果仍有预像——未来之事的圣事——供你们的婚姻来象征；或者，如果现在仍给那道命令留有余地，即\BibleQuote{你们要生育繁殖}；也就是说，如果还没有另一道命令超越它：\BibleQuote{时间已经缩短了；从今以后，那些有妻子的要像没有一样}；因为，当然，通过命令节制并约束交合——我们种族的繁衍，（这后一个命令）已经废除了那\BibleQuote{你们要生育繁殖}。此外，我认为，每一道宣告和安排都是同一天主（的行为）；祂确实在起初，通过给予婚姻结合之缰绳一种宽松的纵容，发出了种族的播种，直到世界被充满，直到新纪律的材料臻于成熟；然而现在，在时代的尽头，祂已经阻止了祂所发出的（命令），并收回了祂所给予的纵容；并非没有合理的理由：起初的扩展和终点的限制。松懈总是被允许于（事物的）开端。一个人种植一片树林并让它生长的原因，是他在自己的时候可以砍伐它。那树林是旧秩序，正被新福音修剪，其中\BibleQuote{斧子已经放在树根上}。同样，\BibleQuote{以眼还眼，以牙还牙}现在已经陈旧，自从\BibleQuote{不可以恶报恶}变得年轻。此外，我认为，即使就\Emph{人的}制度和法令而言，后来者总是胜过早期者。\par

% structure: p0015 source=h3 target=section reason=relative-heading-rank; base=h3

\section{第七章 即使在旧纪律中也不乏强制一夫一妻的先例。但在这点及在其他方面，新约带来了更高的完美。}

此外，为什么我们不应该从（储存的）原始先例中认识到那些在纪律上与后来的（事物秩序）相通，并向新事物传递古风的典型形式？看，在旧律法中，我发现了修枝刀应用于重复婚姻的放纵。\BibleBook{肋未}{纪}中有警告：\BibleQuote{我的司铎不可多婚。} 我甚至可以断言，不是一次而终的就是多。不统一的就是数目。简而言之，统一之后才有数目。此外，统一就是一切一次而终的。但为基督保留的，如同在其他所有点上一样，也有\BibleQuote{法律的成全。} 因此，在\Emph{我们}中间，规定被更充分、更小心地设立，即被选入司祭阶层的人必须是一次婚姻的男子；这条规则被如此严格地遵守，以至于我记得有些因\OriginalTerm{再婚}{digamy}而被撤职。但你会说：\Quote{那么他所除外的所有人，都可以（多次结婚）。} 如果我们认为对司铎不合法的对平信徒合法，我们就枉然了。难道我们平信徒不也是司铎吗？经上记载：\BibleQuote{祂使我们成为国度，成为祂天主和父的司铎。} 是教会的权威，以及通过圣职阶层的共席而获得圣洁的荣誉，建立了圣职阶层与平信徒之间的区别。因此，在没有教会圣职阶层共席的地方，你独自献祭、施洗，并为自己做司铎。但哪里有三个人，哪里就是教会，即使他们是平信徒。因为每个人凭自己的信德生活，在天主面前没有偏袒；因为使徒也说，不是听法律的，而是行法律的，在主面前成义。因此，如果你在自己身上拥有司铎的\Emph{权利}，在必要的情况下，你也有责任拥有司铎的\Emph{纪律}，无论何时有必要拥有司铎的权利。如果你是\OriginalTerm{再婚者}{digamist}，你会施洗吗？如果你是再婚者，你会献祭吗？一个再婚的平信徒充当司铎，比起司铎本人若成为再婚者则被剥夺司铎职权，是何等更大的（罪行）！你说：\Quote{但对必要，给予宽容。} 没有哪种必要是不可避免却可原谅的。总之，避免被发现有再婚的罪责，你就不会使自己面临执行再婚者不可合法执行的事务的必要。天主愿意我们所有人都如此准备，随时随地去承担祂的圣事的（职责）。有\BibleQuote{唯一的天主，唯一的信德}，也有唯一的纪律。事实确实如此，除非平信徒也遵守那些指导挑选长老的规则，否则长老从哪里来呢？他们是从平信徒中被选出来担任此职的。因此，我们必须坚持认为，禁止第二次婚姻的命令\Emph{首先}是针对平信徒的；只要除了平信徒没有其他人能成为长老，只要他曾经是\Emph{一次而终}的丈夫。\par

% structure: p0017 source=h3 target=section reason=relative-heading-rank; base=h3

\section{第八章 即使承认第二次婚姻是合法的，然而并非所有合法之事都有裨益。}

现在我们姑且承认重复婚姻是\OriginalTerm{合法}{lawful}的，如果一切合法的事都是\OriginalTerm{善}{good}的。同一宗徒高呼：\BibleQuote{凡事都可行，但不全有益。}请问，被称为\Quote{不\OriginalTerm{有益}{profitable}}的，能称为善吗？如果连无助于救恩的事也是\Quote{合法}的，那么连不善的事也是\Quote{合法}的了。但你更应选择哪一种呢？是因其\Quote{合法}而为善的，还是因其\Quote{有益}而为善的呢？我认为\Quote{\OriginalTerm{许可}{licence}}与救恩之间存在着巨大差别。关于\Quote{善}，并未说\Quote{它是合法的}，因为\Quote{善}不是期待被允许，而是要被接受。但所谓\Quote{\OriginalTerm{准许}{permission}}是指对其是否为\Quote{善}有疑虑的事；同样也可能不被准许，如果没有某种首要的（外在）原因：——也就是说，第二次婚姻之所以被准许，是\Emph{由于\OriginalTerm{不节制}{incontinence}的危险}：——因为，若非某种（非绝对）善的\Quote{许可}被置于我们的选择之下，就无法证明谁自愿服从天主的旨意，谁服从自己的权力；我们谁遵循现实，谁又抓住许可的机会。\Quote{许可}多半是对纪律的考验；因为纪律是通过考验得到证明，而考验是通过\Quote{许可}来运作的。由此便出现了\BibleQuote{凡事都可行，但不全有益}的情况，只要（仍然真实的是）凡获得\Quote{准许}的人就被考验，并因此在考验过程中（就特定\Quote{准许}的情形）被审判。此外，宗徒们也有\Quote{许可}结婚并带着妻子同行。他们也有\Quote{许可}\Quote{靠福音生活}。但当需要时，他\Quote{没有使用这权利}，这激励我们效法他的榜样，教导我们：我们的考验在于，\Quote{许可}为节制的实践证据奠定了基础。\par

% structure: p0019 source=h3 target=section reason=relative-heading-rank; base=h3

\section{第九章 再婚是奸淫的一种，婚姻本身也受指责，因其近似奸淫}

如果我们深入理解他的意思并加以解释，再婚就不得不被称为一种\OriginalTerm{奸淫}{fornication}。因为，既然他说已婚的人有这样的挂虑，\BibleQuote{如何彼此悦乐}（当然，不是\Emph{道德上}的，因为他不指责好的挂虑）；并且他希望人们理解为他们在挂虑服饰、装饰和各种个人魅力，以增加诱惑力；此外，以美貌和服饰取悦人是血肉\OriginalTerm{贪欲}{concupiscence}的特性，而这是奸淫的原因：请问，再婚是否似乎接近于奸淫，因为其中发现了属于奸淫的成分？主亲自说：\BibleQuote{凡注视妇女，有意贪恋她的，他已在心里奸淫了她。}但注视她有意结婚的人，是做得更轻还是更重呢？如果他甚至娶了她呢？——他若不是有意贪恋她而注视她，就不会娶她；除非可能娶一个你未曾见到或没有欲望的妻子。我承认，已婚男人或未婚男人贪恋另一个女人之间有很大差别。然而，每个女人，甚至对未婚男人来说，只要她属于别人，就是\Quote{别人}；而她成为已婚女人的手段与成为\OriginalTerm{奸妇}{adulteress}的手段并无不同。似乎是法律造成了婚姻与奸淫之间的区别；通过不法性的差异，而不是事物本身的本质。此外，在所有男女中产生婚姻和奸淫的事是什么？当然是肉体的\OriginalTerm{混合}{commixture}；主将这种贪欲与奸淫等同。然后，有人说：\Quote{那么，你连第一次——即单身——婚姻也要摧毁吗？}这并非毫无理由；因为第一次婚姻也包含着奸淫的本质。因此，男人不接触女人是最好的；因此，\OriginalTerm{童贞}{virginity}是首要的圣洁，因为它与奸淫无关。既然这些考虑即使在第一次和唯一婚姻的情况下也能推进\OriginalTerm{节制}{continence}的目标，那么它们对于拒绝再婚岂不更提供了预先的判断？如果天主曾一次性赐予你结婚的放纵，你当感恩。而且，如果你不知道祂再次赐予了你那放纵，你将更是感恩。但如果你无\OriginalTerm{节制}{moderation}地使用它，你就是滥用放纵。\OriginalTerm{节制}{moderation}被理解为源自\Emph{\OriginalTerm{限度}{modus}}，即一个极限。你从最纯洁童贞的最高等级因结婚而跌落还不够，你还要跌入第三、第四、甚至更多，因为你未能在第二阶段保持节制；因为论述缔结第二次婚姻的人甚至不愿禁止更多。因此，让我们每日结婚。在结婚中，让我们被末日赶上，如同索多玛和哈摩辣；那是\Quote{祸哉}宣告给\Quote{怀孕的和哺乳的}的日子，即已婚者和不节制者：因为婚姻带来子宫、乳房和婴孩。结婚何时结束？我相信是在生命结束后！\par

% structure: p0021 source=h3 target=section reason=relative-heading-rank; base=h3

\section{第十章 主题的应用。守寡的益处}

我们当弃绝肉性之事，好能最终结出属神的果实。要抓住这个机会——虽非热切渴望，却是有利的——即不必再向任何人偿还债务，也无人向你偿还！你已不再是债务人。有福的人！你释放了你的债务人；承受这损失吧。倘若你发觉我们所谓的损失其实是收益，那又如何？因为节制将成为一种途径，使你借以获取圣德这一巨大财富；藉着肉身的俭省，你将赢得圣神。让我们省察自己的良心吧，（看看）当一个人偶然失去妻子时，他感到自己有多么不同。他品尝属神的事。若他向主祈祷，他便接近天堂。若他俯首于圣经，他\Quote{全在其中}。若他咏唱圣咏，他自得其乐。若他驱逐魔鬼，他对自己充满信心。因此，宗徒添加了（推荐）暂时禁欲，以增加祈祷的效力，好让我们知道，那\Quote{暂时}有益的事，我们应常行，使之永远有益。每日、每时，祈祷为人所必需；当然，节制亦然，因为祈祷是必需的。祈祷发自良心。若良心羞惭，祈祷也羞惭。是圣神引领祈祷到达天主。若圣神因良心的羞惭而自责，它怎敢引领祈祷到达祭台？因为，若祈祷羞惭，连那神圣的（祈祷）侍奉者自身也满面通红。因为旧约中有一先知预言：\BibleQuote{你们应当是圣的，因为天主是圣的；}又说：\BibleQuote{你与圣者相处，必成为圣者；与无辜的人相处，必成为无辜；与被选的人相处，必成为被选的。}因为我们的责任是在主的训导中行走，正如\Quote{相称}的，而非依照肉体污秽的贪欲。因为宗徒也是如此说：\BibleQuote{随从肉性的思想，导入死亡；但随从圣神的思想，导入永生，在耶稣基督我们的主内。}再者，藉着圣女先知普黎斯卡，福音如此宣讲：\Quote{神圣的侍奉者知道如何侍奉圣德。}\Quote{因为洁净}，她说，\Quote{是和谐的，她们看见异象；并且，她们俯首向下，甚至听见清晰的声音，既是有益的，又是隐秘的。}倘若这种（属神官能的）迟钝——即使在初婚中允许肉性本性的活动——就使圣神退避，那么，当它在再婚中被带入时，又当如何呢？\par

% structure: p0023 source=h3 target=section reason=relative-heading-rank; base=h3

\section{妻子越多，心神越分心}

因为（在此情形下）羞耻是双倍的；在再婚中，两个妻子困扰着同一个丈夫——一个在灵里，一个在肉身上。因为你不能恨第一个妻子，你对她保留着甚至更虔诚的感情，因她已被接到主的面前；你为她的灵魂祈求，为她奉献年度祭品。那么，你将带着如同你在祈祷中纪念的那么多妻子站立在主面前吗？你将献祭为两个吗？你将把那两个（妇女）托付给一位因一夫一妻制而受按立（至神圣职事）或因童贞而被祝圣的司铎的职务，被只结过一次婚的寡妇们环绕吗？你的祭品会毫无愧色地上升吗？并且，在善良心智的其它一切（恩惠）中，你会为自己和你的妻子祈求贞洁吗？\par

% structure: p0025 source=h3 target=section reason=relative-heading-rank; base=h3

\section{为再婚辩护的常见借口。它们的虚妄，尤其在基督徒身上被指出}

我知道那些借口，我们用它们来掩饰我们无法满足的肉身欲望。我们的借口是：需要倚靠的支柱；要管理的房屋；要治理的家庭；要守护的箱子和钥匙；要分配的纺纱；要照料的食物；要通常减轻的忧虑。当然，只有已婚男子的家庭才兴旺！独身者的家庭、阉人的地产、军人或那些不带妻子旅行的人的财产，都已化为废墟！难道我们不也是军人吗？确实是军人，服从于更严格的纪律，因为我们服从于如此伟大的统帅？难道我们不也是这世上的旅客吗？而且，基督徒啊，你为何被如此束缚，以至于没有妻子就不能（这样旅行）？\Quote{在我目前（寡居）的状态，也需要一个家务伴侣。}（那么）娶一位属神的妻子吧。从寡妇中挑选一位在信德上美丽、以贫穷为嫁妆、以年龄为印记的。你将（由此）缔结一桩良缘。多个\Emph{这样的}妻子是悦乐天主的。\Quote{但基督徒关心子嗣}——然而明天不会到来！天主的仆人，已使自己脱离世界，还会渴望继承人吗？如果一个人从第一次婚姻中没有孩子，这能成为他再次结婚的理由吗？那么，他是否以此作为首要的（由此而来的）好处，即他应该渴望更长的寿命，而宗徒本人却急于要\Quote{与主同在？}确实，他将最无牵挂地面对迫害，最坚定地面对殉道，最慷慨地分发财物，最节制地获得；最后，他将无忧无虑地死去，当他留下儿女之后——也许他们还会在他的坟前为他举行最后的仪式！那么，难道是为了国家的远景才缔结这样的（婚姻）吗？恐怕国家会衰亡，如果没有新生代被培养起来？恐怕法律的权利、商业的部门会完全没落？恐怕庙宇被完全遗弃？恐怕没有人发出欢呼：\Quote{狮子给基督徒？}——因为那些寻求子嗣的人渴望听到的就是这些欢呼！让众所周知的孩子负担——尤其是在\Emph{我们的}情况中——足以劝人守寡：孩子们是法律强迫人承担（照管）的；因为没有智慧人会自愿渴望儿子！那么，如果你成功让你的新妻子也充满你自己的良心顾虑，你会怎么做？你会借助药物中断怀孕吗？我认为，对\Emph{我们}来说，伤害正在出生的（孩子）并不比伤害已出生的更合法。但也许在你妻子怀孕的那个时候，你会有勇气向天主祈求解救如此严重的忧虑，而这曾在你自己的能力范围内时，你却拒绝了？我猜，你预谋寻找的目标会是一些（天生）不孕的女人，或是一些年龄已感年岁寒意的女人。这是相当谨慎的做法，尤其是值得一位信徒所做的！因为当天主愿意时，没有哪个女人是我们相信曾在不能生育或年老时怀上（孩子）的！如果任何人因自己这种远见的僭越，激起天主的竞争，他更可能这样做。最后，我们知道在弟兄中有一个案例，其中一人为了女儿的缘故娶了一个不能生育的女人作为第二次婚姻，并且既第二次做了父亲，也第二次做了丈夫。\par

% structure: p0027 source=h3 target=section reason=relative-heading-rank; base=h3

\section{第十三章 来自外邦人和教会中的榜样，以加强上述劝勉}

在我这番劝勉中，最亲爱的弟兄，甚至外教人的例子也被添加进来；这些例子我们也（像别人一样）常引以为证，当任何善事和悦乐天主的事，即使在\Quote{外邦人}中，也被承认并受到表扬的时候。简言之，一夫一妻在外教人中如此备受尊崇，以至于合法出嫁的处女，常有一位从未结过婚的女子被指定为她们的伴娘；如果\Emph{你说}\InnerQuote{这是为了预兆的缘故}，当然是为了一个\Emph{美好的}预兆；再者，在某些庄严的礼仪和公务中，独婚享有优先地位；无论如何，一位神祇祭司的妻子必须只结一次婚，这也是（神祇祭司本人）的律法。大司祭本人不得再婚的事实，当然是一夫一妻的荣耀。然而，当撒殚摹仿天主的圣事时，这对我们是一种挑战；不，更是一种羞愧，如果我们迟迟不肯向天主展示一种节制，而有些人却通过有时是永久的贞女生活、有时是永久的寡居生活，将这节制献给魔鬼。我们听说过维斯塔的贞女，阿哈雅城的朱诺（贞女），德尔斐人中的阿波罗（贞女），以及某些地方的米内瓦和狄安娜（贞女）。我们也听说过节制的\Emph{男子}，（其中）有著名的埃及公牛之司祭；还有妇女们，（奉献）给非洲的刻瑞斯，她们为了尊敬她甚至自愿放弃婚姻，如此活到老年，从此回避与男性的一切接触，甚至儿子的亲吻。的确，魔鬼在肉欲之外，还发现了一种引人丧亡的贞洁；这样，拒绝那助人得救的贞洁的基督徒，其罪过就更加深重！外教人中一些因坚守独婚而闻名妇女，也将为我们作证：比如狄多，她身为流落异乡的难民，本应无需外界催促就渴望与国王结婚，却因害怕经历第二次结合，反而宁愿\Quote{焚烧}也不愿\Quote{结婚}；又比如著名的卢克雷蒂娅，尽管她只是被迫、违背自己意愿地遭受了一次陌生男子的侵犯，仍用自己鲜血洗净被玷污的肉体，免得自己活着，而在她自己眼中已不再是独婚者！稍加留意，你就能从我们自己的（姐妹）中找到更多例子；而\Emph{那些}例子，确实优于前者，因为活着保持贞洁比为贞洁而死更为伟大。失去福乐而献出生命，比为那福乐宁愿彻底死去却活着保存它更容易。因此，有多少男子和多少女子，在教会圣秩中，因节制而获得其地位，他们宁愿与天主结婚；他们恢复了他们肉身的尊荣，并已将自己奉献为那（未来）时代的儿子，在他们自身内杀死了私欲以及那整个（倾向），这倾向是不能进入乐园的！由此可以推测，那些希望被接纳入乐园的人，终归应当开始从那使乐园完好无损的东西中脱离出来。\par

\SourceRecord{Translated by S. Thelwall.；Ante-Nicene Fathers；Vol. 4.；Edited by Alexander Roberts, James Donaldson, and A. Cleveland Coxe.；Buffalo, NY: Christian Literature Publishing Co.,；1885.；<http://www.newadvent.org/fathers/0405.htm>.}
